% !PRINTED 63
\weberpage{63}
%<<CONT

\begin{equation}
 n\Phi(x_1,x_2,\ldots,x_m)
 =x_1\Phi'(x_1)+x_2\Phi'(x_2)+\cdots+x_m\Phi'(x_m),
\tag{5}
\end{equation}
of which formula (15) of the preceding paragraph is a special case, and for $\nu=2$:
\begin{equation}
 n(n-1)\Phi(x_1,x_2,\ldots,x_m)
 =
 \sum_{i,k=1}^{m}x_i x_k\Phi''(x_i,x_k),
\tag{6}
\end{equation}
where the sum is to be extended from $i=1$ to $i=m$ and from $k=1$ to $k=m$, so that every term with unequal $i,k$ occurs twice in the sum.

If, when the $\alpha$ are subject to the condition (4), we put
\begin{equation}
 \Phi_\nu(\xi,x)=
 \sum {\xi_1^{\alpha_1}\xi_2^{\alpha_2}\cdots\xi_m^{\alpha_m}
 \over \Pi(\alpha_1)\Pi(\alpha_2)\cdots\Pi(\alpha_m)}
 D_{\alpha_1,\alpha_2,\ldots,\alpha_m}\Phi,
\tag{7}
\end{equation}
then by (4) of \S~16,
\begin{equation}
 \Phi(x+\xi)=\Phi(x)+\Phi_1(x,\xi)+\Phi_2(x,\xi)+\cdots+\Phi_n(x,\xi),
\tag{8}
\end{equation}
and since the left-hand side remains unchanged when $x$ and $\xi$ are interchanged, the relation follows:
\begin{equation}
 \Phi_{n-\nu}(x,\xi)=\Phi_\nu(\xi,x),
\tag{9}
\end{equation}
and therefore in particular
\begin{equation}
 \Phi_n(x,\xi)=\Phi(\xi).
\tag{10}
\end{equation}
The function $\Phi_\nu(x,\xi)$, regarded as a function of $x$, is called the \gesperrt{$\nu$th polar} of the function $\Phi$ for the system of values $\xi$.

We shall now specialize formula (3) to the case of a binary form ($m=2$). We denote the variables by $x,y$ and put, for abbreviation,
\[
 \Phi(x,y)=u,\qquad D_{h,\nu-h}\Phi=u_h.
\]
% [source-original sic] Printed p.63 cites (4), although this specialization is derived from (3); corrected in translation.
From (3) we obtain
\begin{equation}
 {\Pi(n)\over\Pi(n-\nu)}u
 =
 \sum_{h=0}^{\nu}{\Pi(\nu)\over\Pi(h)\Pi(\nu-h)}
 u_h x^h y^{\nu-h},
\tag{9}% [source-original sic] Printed p.63 repeats (9) after (9) and (10); retained verbatim.
\end{equation}
where $\nu$ can assume any value not greater than $n$.
